\section{Client Resource Constraints}

NAS is computationally demanding even on high-end hardware, and individual experiments can require extensive GPU resources for weeks~\cite{TODO}. FL clients often have much less compute and memory available. Research therefore focuses on reducing resource use. The challenges below distinguish whether a client can execute its assigned workload, the total cost of the search and its communication overhead.

\subsubsection{Challenge 1: Client Capacity Limits the Search Work Assigned Per Round}

NAS-in-FL becomes infeasible when an assigned search workload exceeds a client's compute or memory. Hardware differences compound this problem because the same workload may fit one client but not another~\cite{fedoras_2022,fednas_over_hetero_mobile_devices_2022,dfnas_2020,peaches_2024,supernet_trade-off_fednas_2025}. Differentiable supernet-based strategies are particularly demanding because clients update several candidate operations within a shared supernet~\cite{fednas_2021,dp-fnas_2020,feathers_2023,fed_auto_mri_2023}. Other strategies usually assign one candidate at a time, but even that candidate can exceed a low-end client's capacity. This challenge concerns whether clients can execute their assigned task. Challenge 9 concerns how long capable clients take to complete it.

\noindent\textbf{Solution 1a: Use resource-efficient search components.} Practitioners can restrict the search space to inexpensive components, such as depthwise-separable convolutions and mobile-oriented CNN blocks~\cite{fed_auto_mri_2023,intrusion_detection_fednas_2024}. This makes every candidate cheaper to train, but may exclude costly components that improve prediction accuracy. Solution 10a instead tests candidate modules on server data.

\noindent\textbf{Solution 1b: Send clients subnets tailored to their capacity.} Instead of distributing the complete supernet, the server can send each client a subnet that fits its compute and memory limits. For clients with fewer resources, the server can reduce the subnet's depth or width, or assign fewer local optimization steps~\cite{dfnas_2020,fedoras_2022,apons_2023,divide-and-conquer_fednas_2023,peaches_2024,fgfl_2024,supernet_trade-off_fednas_2025,optimize_fednas_edge_2021,efnas_2024,load_monitoring_fednas_2025,perfedrlnas_2024,rafl_2024}. As a result, some operations or widths may receive fewer updates, as discussed in Challenge 6.

\noindent\textbf{Solution 1c: Search trainable adaptation modules over a frozen backbone.} When a suitable pretrained model is available, clients can keep its backbone fixed and search for a small combination of lightweight adaptation modules. Only these modules and their architecture parameters require local training and exchange, reducing computation and update size~\cite{fednaspet_2023}. The benefit depends on how well the fixed backbone transfers to the clients' tasks.

\noindent\textbf{Cross-cutting solution 7d.} Architecture-agnostic distillation lets small client models teach a server-held supernet without requiring clients to run it~\cite{hapfnas_2025}.

\subsubsection{Challenge 2: Architecture Search Adds Computational Cost}

Even when clients can handle each round, the complete architecture search can take too long. Black-box strategies may train several candidates in separate federated evaluations~\cite{decnas_2022,rt_fed_evo_nas_2020}, while weight-sharing strategies repeatedly update a large supernet or sampled subnets~\cite{fednas_2021,superfednas_2024}. Split-model execution can also leave clients idle while other segments run~\cite{nasfly_2024}. Some strategies discard the search-trained weights and retrain the selected architecture from a new initialization, adding more local epochs and aggregation rounds~\cite{medical_image_segmentation_fednas_2024,privacy_preserving_fednas_for_edge_2025}. This challenge arises from repeated candidate evaluation, idle resources or search work that produces no reusable weights. Challenge 3 addresses communication volume.

\noindent\textbf{Solution 2a: Use multi-fidelity evaluation for candidate architectures.} Black-box strategies can first evaluate each candidate with only a few communication rounds. Intermediate validation performance determines which candidates stop early and which receive more rounds. This multi-fidelity approach is used across several black-box search strategies and follows established AutoML methods~\cite{automl_book_hpo_2019}.

For example, \cite{decnas_2022,fednas_over_hetero_mobile_devices_2022} report that the best candidate usually becomes distinguishable within the first few rounds. Their procedures drop weaker candidates and increase the budget in later search iterations. Some evolutionary strategies likewise use a small fixed budget or early stopping~\cite{multi_objective_evo_fl_2020,gpt4-boosted_pso_fednas_2024}.

\noindent\textbf{Solution 2b: Prune a supernet progressively.} Supernet-based strategies can periodically remove weak operations, leaving fewer alternatives to store and train in later rounds~\cite{dfnas_2020,peaches_2024}. Unlike changing the evaluation budget, pruning reduces the search space itself.

Whether the final model needs separate training depends on how its weights are retained, as addressed by Solutions 2e and 2g. Pruning saves more work as search progresses, but early rounds still train the complete supernet and premature pruning can remove useful operations.

\noindent\textbf{Solution 2c: Predict candidate performance.} Performance prediction reduces costly configuration evaluations in AutoML~\cite{automl_book_hpo_2019}. NAS-in-FL strategies can evaluate a small set of architectures and train a predictor to estimate the quality of the rest~\cite{hapfnas_2025,energy_fednas_2025}.

Because private, non-IID data can produce different rankings at each client, a separate predictor can be trained for each client~\cite{hapfnas_2025}. Predictors can also be learned collaboratively without transmitting locally evaluated architectures~\cite{federated_bayesian_optimization_nas_2023}. This reduces client evaluations, but works only if the evaluated architectures represent the relevant search-space regions.

\noindent\textbf{Solution 2d: Evaluate independent candidates or subnets in parallel.} Black-box strategies can train and evaluate different candidates concurrently on separate client groups~\cite{decnas_2022,rt_fed_evo_nas_2020,energy_fednas_2025}.

Supernet-based strategies can sample several subnets, assign each to a client group and merge their updates into the shared supernet~\cite{finch_2024}. This can reduce search time when enough clients are available, but different group data can bias candidate comparisons~\cite{decnas_2022}. Challenge 6 addresses this problem.

\noindent\textbf{Solution 2e: Reuse search-trained weights for candidate evaluation and deployment.} In weight-sharing strategies, candidates can inherit the operation weights associated with their architecture~\cite{fl-agnns_2023,optimize_fednas_edge_2021}. This avoids training every candidate from scratch and can initialize the selected subnet before limited fine-tuning~\cite{medical_image_segmentation_fednas_2024}. Here, architecture selection remains separate from supernet training. Solution 2g instead produces the architecture and deployable weights together. Shared weights may still rank candidates poorly or provide a weak final initialization.

\noindent\textbf{Solution 2f: Use split-training idle periods for local search.} Temporary input and output components can make an assigned model segment executable on its own. During idle periods, a client can then sample and train alternative branches without waiting for the other model segments~\cite{nasfly_2024}. The assigned segment can also provide the feature-based teaching signal described in Solution 7d. This converts waiting time into search work.

\noindent\textbf{Solution 2g: Jointly derive an architecture and its deployable weights.} One-stage supernet strategies can train model weights while narrowing the architecture choices until the search produces a discrete model with trained weights~\cite{dfnas_2020,fed_meta_nas_2025}. Reinforcement-learning strategies over a shared supernet can similarly use later rounds to fine-tune the final model as sampling concentrates on one path~\cite{perfedrlnas_2024}. Unlike Solution 2b, this makes final model production part of search. The selected path may remain undertrained if competing paths received most updates.

\subsubsection{Challenge 3: Architecture Search Increases Communication Costs}

FL with a predefined architecture exchanges one model's parameters. NAS-in-FL may send larger updates, synchronize more often and exchange many candidate models. A supernet stores weights for alternative operations and can be much larger than the selected architecture. Differentiable strategies may also send architecture parameters throughout search~\cite{fedoras_2022,hapfnas_2025,fednas_2021,fedregnas_2025}.

\noindent\textbf{Solution 3a: Use compact codes to identify known candidates.} When clients already hold a compatible master model or trained supernet, the server can send a compact selector or genotype instead of the candidate's parameters. Clients reconstruct the subnet locally and return its evaluation or updated weights~\cite{rt_fed_evo_nas_2020,medical_image_segmentation_fednas_2024,intrusion_detection_fednas_2024}. Unlike Solution 7b, this requires a shared search space and common retained weights.

\noindent\textbf{Solution 3b: Synchronize architecture parameters less frequently.} Clients can update and upload architecture parameters less often than model weights. One approach transmits them every \(H_{\alpha}\) rounds and stops once the architecture is fixed, while weight training continues~\cite{fedregnas_2025}.

\noindent\textbf{Solution 3c: Aggregate search updates hierarchically.} A nearby aggregator can combine updates from several clients before forwarding one result to the remote server. Most aggregation then uses local links, reducing transfers over the wider network~\cite{finch_2024}. In one non-IID CIFAR-10 evaluation, this reduced traffic at 80\% target accuracy from 3.29--4.68 GB to 1.64 GB~\cite{finch_2024}.

The selected architecture also determines later communication costs. A model chosen only for predictive performance or local computation may have many parameters that FL must transmit in every training round. It can therefore meet compute and memory limits while exceeding the transfer budget~\cite{multi_objective_evo_fl_2020}. This challenge covers communication during search and during later training of the selected model.

\noindent\textbf{Applicable cross-cutting solution.} Solution 4b can include model size or expected transmitted bytes in candidate scores, allowing the search to consider later communication costs before selecting an architecture~\cite{multi_objective_evo_fl_2020,fedregnas_2025}.
